\documentclass[11pt]{article}
\usepackage[margin=1in]{geometry}
\usepackage{booktabs}
\usepackage{longtable}
\usepackage{amsmath}
\title{Supplementary Information: Active SolvAI}
\date{}
\begin{document}
\maketitle
\section{Scope and target convention}
The simulation-reconstruction target and experimental hydration-free-energy target are distinct. Raw Arbalest annihilation $dH/d\lambda$ is retained with its native sign; a hydration free-energy integral uses the documented sign conversion.
\section{Identity and exposure audit}
The Phase 1 cohort contains the 72 ARROW molecules for which all three PIMD2 probe windows completed. Molecule identity is fixed by the parent benchmark InChIKey and connectivity key; no molecule was removed according to its response or endpoint error. These molecules are development-exposed through the parent SolvAI evaluation and are not an external validation cohort.
\section{Parent SolvAI reproducibility}
Parent metrics, model hashes and smoke predictions are regenerated by \texttt{scripts/reproduce\_parent.py}.
\section{Response inventory}
The three campaigns attempted 76 molecules at each of nominal $\lambda=0.1$, 0.5 and 0.9. Successful counts were 72, 74 and 73, respectively; 72 molecules completed all three. Each successful window used two beads, 500 equilibration steps and 2,500 production steps at 0.002 ps per step. Sequential prefix summaries used the first 10\%, 20\%, 40\%, 70\% and 100\% of stored frames.
\section{Phase 1 model matrix}
All one-, two- and three-point subsets of nominal $\lambda=\{0.1,0.5,0.9\}$ are predeclared. Raw actual observations, actual-minus-predicted residuals, shuffled observations and simple generic controls are compared on identical molecules and outer splits.
The primary correction was a standardized ridge model whose regularization parameter was selected only within the outer training data. It used the total SYSTEM $dH/d\lambda$ residual at all three points. The matched frozen SolvAI baseline was reconstructed with the same 1,280 public labels, outer-training ARROW labels, sample weights, corrected response teachers and three-member ExtraTrees endpoint ensemble as the parent confirmation.

Across five partitions the frozen baseline MAE was 0.1864 and the residual-corrected MAE was 0.1899 kcal mol$^{-1}$. The candidate-minus-baseline paired change was $+0.00348$ (95\% bootstrap interval $+0.00096$ to $+0.00609$). The matched predicted-response, raw-response and mean shuffled-residual MAEs were 0.1916, 0.1907 and 0.1893 kcal mol$^{-1}$, respectively.

\section{Sparse posterior diagnostic}
Generic and molecule-conditioned Gaussian models were fitted to the three response coordinates with Ledoit--Wolf covariance shrinkage. Held-out observations were conditioned using their five-contiguous-block standard errors. Molecule-conditioned interpolation improved several one- and two-point hidden-coordinate tests, but hidden-point MAE remained at least 2.4 kcal mol$^{-1}$ and the intervals were wide. This is not a dense-curve validation and cannot establish accurate thermodynamic integration.
\section{Negative results}
All killed directions are recorded in \texttt{NULL\_RESULTS.md} and the append-only experiment ledger.
\end{document}
